% Generated draft: requires mathematical and source review.
\documentclass[11pt]{article}
\usepackage[margin=25mm]{geometry}
\usepackage{fontspec}
\setmainfont{Latin Modern Roman}
\usepackage{amsmath,amssymb,mathtools,unicode-math}
\setmathfont{Latin Modern Math}
\usepackage{enumitem,needspace,xurl,hyperref}
\setlength{\parindent}{0pt}
\setlength{\parskip}{5pt}
\setlength{\emergencystretch}{3em}
\newcommand{\N}{\mathbb{N}}
\newcommand{\Z}{\mathbb{Z}}
\newcommand{\Q}{\mathbb{Q}}
\newcommand{\R}{\mathbb{R}}
\newcommand{\C}{\mathbb{C}}
\newcommand{\F}{\mathbb{F}}
\newcommand{\A}{\mathbb{A}}
\newcommand{\PP}{\mathbb{P}}
\newcommand{\E}{\mathbb{E}}
\newcommand{\eps}{\varepsilon}
\newcommand{\dd}{\,\mathrm{d}}
\begin{document}
\section*{UnsolvedMath ID: 701696}
\textbf{Source problem code:} catalogue entry
\subsection{Definitions and mathematical statement}
For a finite group \(X\), let \(\operatorname{Irr}_{\C}(X)\) be the set of irreducible complex characters of \(X\), and define its set of character degrees by
\[
\operatorname{cd}(X)=\{\chi(1):\chi\in\operatorname{Irr}_{\C}(X)\}.
\]
This is a set, so multiplicities of degrees are disregarded. Here an almost simple group of Lie type means a group \(H\) for which there is a finite nonabelian simple group \(S\) of Lie type such that
\[
S\leq H\leq\operatorname{Aut}(S),
\]
where \(S\) is identified with its group of inner automorphisms.

The question attributed to A. Iranmanesh is whether, for every such \(H\) and every finite group \(G\),
\[
\operatorname{cd}(G)=\operatorname{cd}(H)
\quad\Longrightarrow\quad
\exists\,A\trianglelefteq G\ \text{such that }A\text{ is abelian and }G/A\cong H.
\]

The supplied record includes a \emph{Comment of 2025} reporting that the conjecture was confirmed for projective general linear and projective general unitary groups of dimension \(3\), citing F. Shirjian and A. Iranmanesh (2020). The comment does not specify the field sizes or exceptional parameters covered by that result; those details are left unresolved here.
\subsection{Short English statement}
If a finite group has exactly the same set of irreducible complex character degrees as a given almost simple group of Lie type, must it have an abelian normal subgroup whose quotient is isomorphic to that almost simple group? The supplied record labels the problem open and reports a dimension-\(3\) linear and unitary special case.
\subsection{Sources}
\begin{itemize}
\item Local UnsolvedMath registry, permanent ID \(701696\), exact source code \texttt{catalogue entry}, supplied complete source record; upstream revision \texttt{d371ee917ac9647fbb06cd6fc529deec6f3dda82}.
\item A. Iranmanesh, Problem 20.38 in \emph{The Kourovka Notebook (Unsolved Problems in Group Theory)}, 20th issue (2022); the supplied background also identifies the 21st issue edition (2026), arXiv:1401.0300v48. The registry describes the collection as edited by E. I. Khukhro and V. D. Mazurov.
\item F. Shirjian and A. Iranmanesh, \emph{Illinois Journal of Mathematics} \textbf{64}, no. 1 (2020), 49--69, cited in the supplied \emph{Comment of 2025} for the dimension-\(3\) projective general linear and unitary cases. No article title or link is supplied. This reported result and the registry's open label are supplied source assessments, not independently verified here.
\end{itemize}
\end{document}
